\documentclass[11pt]{article}
\usepackage[margin=1in]{geometry}
\usepackage{amsmath,amssymb,amsthm}
\usepackage{booktabs,array}
\usepackage{enumitem}
\usepackage{hyperref}
\usepackage{xcolor}
\usepackage{graphicx}
\usepackage{subcaption}

\title{Experiment Design: Analysis of Muon with Nesterov Momentum}
\author{}
\date{}

\newcolumntype{L}[1]{>{\raggedright\arraybackslash}p{#1}}

\begin{document}
\maketitle

\section{Overview}

This document describes the experiment plan for the paper ``Analysis of Muon with Nesterov Momentum.'' First implement all candidate methods, then use ablations to choose the randomized configuration, and finally compare the chosen method against strong baselines.

\begin{table}[ht]
\centering
\begin{tabular}{L{0.15\linewidth}L{0.43\linewidth}L{0.30\linewidth}}
\toprule
\textbf{Exp.} & \textbf{Purpose} & \textbf{Output} \\
\midrule
E1 & Ablate inexact solvers & solver selected \\
E2 & Sweep normalized rank $k/d$ & Selected rank \\
E3 & Sweep NS steps $q$ & Selected $q$ \\
E4 & Sweep batch size $B$ & Selected batch size \\
E5 & Compare the final randomized method against baselines & Final table and training curves \\
E6 (optional) & Explore sample complexity via $\varepsilon$-threshold stopping & Empirical $K$ versus $\varepsilon$ trend \\
\bottomrule
\end{tabular}
\end{table}

\noindent\textbf{Execution order:} implement all candidate methods $\to$ E1 $\to$ E2 $\to$ E3 $\to$ E4 $\to$ E5. Run E6 only after the main comparison is stable.

\paragraph{Shared defaults before selection.} Unless an experiment is actively sweeping a hyperparameter, keep the randomized Muon setup fixed at Nesterov momentum, oversampling $p=10$, one power iteration $h=1$, provisional NS depth $q=5$, provisional normalized rank $k = d/8$, and provisional CIFAR-10 batch size $B=2000$ and Nanogpt batch size $B=1024$. The provisional values for $k$, $q$, and $B$ are only working defaults for earlier ablations; later experiments replace them with the selected values.

\paragraph{Evaluation metric definition.} We use Validation Perplexity for defining the accuracy for Nanogpt and accuracy for the CIFAR-10.


%%%%%%%%%%%%%%%%%%%%%%%%%%%%%%%%%%%%%%%%%%%%%%%
\section{Experiment E1: Inexact Solver Ablation}
%%%%%%%%%%%%%%%%%%%%%%%%%%%%%%%%%%%%%%%%%%%%%%%

\subsection{Goal}

Compare candidate inexact solvers and choose the solver that gives the best accuracy.

\subsection{Protocol}

\begin{enumerate}[nosep]
\item Fix a shared model, training budget, data split, oversampling $p=10$, power iteration $h=1$, provisional rank $k = d/8$, default NS depth $q=5$, and batch size $B=2000$ for CIFAR-10 or $B=1024$ for Nanogpt.
\item Keep Nesterov momentum fixed so that the only solver-side change is the inexact orthogonalization method.
\item Compare four inexact solvers:
\begin{enumerate}[nosep]
    \item Cubic Newton-Schulz: $f(x) = \tfrac{3}{2}x - \tfrac{1}{2}x^3$
    \item Quintic Newton-Schulz with theoretical coefficients: $f(x) = 2x - \tfrac{3}{2}x^3 + \tfrac{1}{2}x^5$
    \item Quintic Newton-Schulz with empirical coefficients: $(a,b,c) = (3.4445,\; -4.7750,\; 2.0315)$
    \item Polar Express (degree-5)
\end{enumerate}
\item For each solver, record training loss, validation/test accuracy, runtime.
\end{enumerate}

\subsection{Table and Figure}
\begin{itemize}[nosep]
\item one validation loss-versus-iteration figure with one line per solver.
\item one final accuracy table with one row per solver.
\end{itemize}

%%%%%%%%%%%%%%%%%%%%%%%%%%%%%%%%%%%%%%%%%%%%%%%
\section{Experiment E2: Rank Sweep}
%%%%%%%%%%%%%%%%%%%%%%%%%%%%%%%%%%%%%%%%%%%%%%%

\subsection{Goal}

Choose the rank $k$ that gives the best performance once the solver has been fixed.

\subsection{Protocol}

\begin{enumerate}[nosep]
\item Fix the solver selected in E1 and keep Nesterov momentum, oversampling $p=10$, power iteration $h=1$, default NS depth $q=5$, and batch size $B=2000$ for CIFAR-10 or $B=1024$ for Nanogpt unchanged.
\item Sweep ranks $k \in \{d/16,\; d/8,\; d/4,\; d/2\}$.
\item For each rank, record training loss, validation/test accuracy, runtime.
\end{enumerate}

\subsection{Outputs}
\begin{itemize}[nosep]
\item one accuracy-versus-iteration figure with one line per rank.
\item one final accuracy table with one row per rank.
\item one singular value decay plot with one line per rank on the last convolution layer of the ConvGroup in CIFAR-10 or the Q,K,V matrices in Nanogpt. (Similar to Figure 5 in the paper: "Low-rank orthogonalization for large-scale matrix optimization with applications to foundation model training".) This is to check the momentum matrix is approximated well enough by the rank-$k$ approximation. We want to say the singular value decay fast.
\end{itemize}

%%%%%%%%%%%%%%%%%%%%%%%%%%%%%%%%%%%%%%%%%%%%%%%
\section{Experiment E3: NS-Step Sweep}
%%%%%%%%%%%%%%%%%%%%%%%%%%%%%%%%%%%%%%%%%%%%%%%

\subsection{Goal}

Choose the NS depth $q$ that gives the best performance after the solver and rank have been fixed.

\subsection{Protocol}

\begin{enumerate}[nosep]
\item Fix the solver selected in E1, the rank selected in E2, Nesterov momentum, oversampling $p=10$, power iteration $h=1$, and batch size $B=2000$.
\item Sweep $q \in \{1, 3, 5, 7, 9\}$.
\item For each $q$, record training loss, validation/test accuracy, runtime.
\end{enumerate}

\subsection{Outputs}

\begin{itemize}[nosep]
\item one accuracy-versus-iteration figure with one line per NS depth.
\item one final accuracy table with one row per NS depth.
\end{itemize}


%%%%%%%%%%%%%%%%%%%%%%%%%%%%%%%%%%%%%%%%%%%%%%%
\section{Experiment E4: Batch-Size Sweep}
%%%%%%%%%%%%%%%%%%%%%%%%%%%%%%%%%%%%%%%%%%%%%%%

\subsection{Goal}

Choose a batch size $B$ after the solver, rank, and NS depth have been fixed.

\subsection{Protocol}

\begin{enumerate}[nosep]
\item Fix the solver selected in E1, the rank selected in E2, the NS depth selected in E3, Nesterov momentum, oversampling $p=10$, and power iteration $h=1$.
\item Sweep batch sizes $B \in \{500,\; 1000,\; 2000, 3000, 4000\}$ for CIFAR-10 and $B \in \{256,\; 512,\; 1024, 2048, 4096\}$ for Nanogpt.
\item For each batch size, record end-to-end accuracy, runtime, throughput, and stability.
\end{enumerate}

\subsection{Outputs}

\begin{itemize}[nosep]
\item one accuracy-versus-iteration figure with one line per batch size.
\item one final accuracy table with one row per batch size.
\end{itemize}


%%%%%%%%%%%%%%%%%%%%%%%%%%%%%%%%%%%%%%%%%%%%%%%
\section{Experiment E5: Comparison with External Methods}
%%%%%%%%%%%%%%%%%%%%%%%%%%%%%%%%%%%%%%%%%%%%%%%

\subsection{Configuration of Our Method}

Before running E5, determine all hyperparameters from the preceding ablations:

\begin{table}[ht]
\centering
\begin{tabular}{ll}
\toprule
\textbf{Parameter} & \textbf{Determined by} \\
\midrule
Orthogonalization solver & E1 \\
Normalized rank $k/d$ & E2 \\
$q$ (NS steps) & E3 \\
Batch size $B$ & E4 \\
Oversampling $p$ & Fixed ($p=10$) \\
Power iteration $h$ & Fixed ($h=1$) \\
$\beta$ (momentum) & Bayesian search \\
Momentum type & Fixed (Nesterov) \\
Learning rate $\eta$ & Bayesian search \\
\bottomrule
\end{tabular}
\end{table}

\subsection{Table and Figure}
\textbf{Summary table:}

\begin{table}[ht]
\centering
\begin{tabular}{lcc}
\toprule
\textbf{Method} & \textbf{Accuracy (mean $\pm$ std)} & \textbf{Time (sec)} \\
\midrule
SGD + Nesterov & & \\
AdamW & & \\
Muon (Polyak, full NS) & & \\
Muon (Nesterov, full NS) & & \\
Muon (Polyak, randomized NS) & & \\
Muon (Nesterov, randomized NS) & & \\
\bottomrule
\end{tabular}
\end{table}

\textbf{Figure:} 
\begin{itemize}[nosep]
    \item one accuracy-versus-iteration figure, one line per method.
    \item one final accuracy table with one row per method.
\end{itemize}


\section{Experiment E6: Sample Complexity}

\subsection{Goal}

Explore the empirical relation between a target accuracy threshold $\varepsilon$ and the number of steps needed to reach it.

\subsection{Protocol}

\begin{enumerate}[nosep]
\item Fix oversampling $p=10$, power iteration $h=1$, the solver selected in E1, the rank selected in E2, the NS depth selected in E3, and the batch size selected in E4. But we need to preselect the max number of iterations $K_{max}$ for choosing $\eta$ to match the theory such as $\eta = K_{max}^{-1/2}$.
\item Choose thresholds $\varepsilon \in \{1e-1, 1e-2, 1e-3, 1e-4, 1e-5\}$.
\item For each $\varepsilon$, train until the monitored gradient norm first satisfies $\|g_k\| \leq \varepsilon$. We also want to run more trails for the same $\varepsilon$ to get a more stable result. And then use the mean of the stopping steps to plot.
\item Record the stopping step $K$ or $K_{max}$ and compute the mean of the stopping steps.
\item Plot the mean of the stopping steps versus $\varepsilon$ and check whether the empirical trend is consistent with the order suggested in A.4 Sample Complexity (Inexact Case) .
\end{enumerate}

\subsection{Output}

\begin{itemize}[nosep]
\item one plot of stopping step $K$ versus $\varepsilon$;
\end{itemize}


%%%%%%%%%%%%%%%%%%%%%%%%%%%%%%%%%%%%%%%%%%%%%%%
\section{CIFAR-10 Results}
%%%%%%%%%%%%%%%%%%%%%%%%%%%%%%%%%%%%%%%%%%%%%%%

All curves are averaged over 50 trials; shaded bands show $\pm 1$ standard deviation. Test accuracies are reported as mean $\pm$ std over trials.

\paragraph{Fixed parameters across experiments.} Each ablation sweeps one parameter while fixing the others; later experiments use the values selected by earlier ones.

\begin{table}[ht]
\centering
\begin{tabular}{lcccc}
\toprule
\textbf{Exp.} & \textbf{Solver} & \textbf{Rank $k$} & \textbf{NS iterations $q$} & \textbf{Batch size $B$} \\
\midrule
E1 & sweep                 & 256         & 5           & 2000 \\
E2 & Quintic (theoretical) & sweep       & 5           & 2000 \\
E3 & Quintic (theoretical) & 256         & sweep       & 2000 \\
E4 & Quintic (theoretical) & 256         & 9           & sweep        \\
E5 & Quintic (theoretical) & 256         & 9           & 1000         \\
\bottomrule
\end{tabular}
\caption{Configuration used in each experiment. ``sweep'' marks the parameter varied in that row; ``prov.'' marks a provisional value used before the corresponding ablation has been run.}
\end{table}

\newcommand{\resultdir}{../muon-randomized-svd/cifar10/experiment_result}

\subsection{E5: Comparison with External Methods}

\begin{figure}[ht]
\centering
\begin{subfigure}[b]{0.49\linewidth}
\includegraphics[width=\linewidth]{\resultdir/E5_val_loss.png}
\caption{Validation loss}
\end{subfigure}
\begin{subfigure}[b]{0.49\linewidth}
\includegraphics[width=\linewidth]{\resultdir/E5_train_loss.png}
\caption{Training loss}
\end{subfigure}
\caption{E5: validation and training loss for each method.}
\end{figure}

\begin{table}[ht]
\centering
\begin{tabular}{lc}
\toprule
\textbf{Method} & \textbf{Test accuracy (mean $\pm$ std)} \\
\midrule
AdamW                         & $0.9106 \pm 0.0026$ \\
SGD-Nesterov                  & $0.9091 \pm 0.0035$ \\
Muon Polyak (full SVD)        & $0.9380 \pm 0.0013$ \\
Muon Polyak (rand SVD)        & $0.9378 \pm 0.0012$ \\
Muon Nesterov (full SVD)      & $0.9377 \pm 0.0011$ \\
\textbf{Muon Nesterov (rand SVD)}      & $\mathbf{0.9383 \pm 0.0014}$ \\
\bottomrule
\end{tabular}
\caption{E5: final TTA test accuracy.}
\end{table}

\subsection{E1: Solver Ablation}

\begin{figure}[ht]
\centering
\begin{subfigure}[b]{0.49\linewidth}
\includegraphics[width=\linewidth]{\resultdir/Ablation_E1_solvers_val_loss.png}
\caption{Validation loss}
\end{subfigure}
\begin{subfigure}[b]{0.49\linewidth}
\includegraphics[width=\linewidth]{\resultdir/Ablation_E1_solvers_train_loss.png}
\caption{Training loss}
\end{subfigure}
\caption{E1: validation and training loss by solver.}
\end{figure}

\begin{table}[ht]
\centering
\begin{tabular}{lc}
\toprule
\textbf{Solver} & \textbf{Test accuracy (mean $\pm$ std)} \\
\midrule
Polar Express          & $0.9384 \pm 0.0014$ \\
Cubic (theoretical)    & $0.9383 \pm 0.0012$ \\
\textbf{Quintic (theoretical)}  & $\mathbf{0.9384 \pm 0.0016}$ \\
Quintic (empirical)    & $0.9382 \pm 0.0013$ \\
\bottomrule
\end{tabular}
\caption{E1: final TTA test accuracy by solver.}
\end{table}

\subsection{E2: Rank Sweep}

\begin{figure}[ht]
\centering
\begin{subfigure}[b]{0.49\linewidth}
\includegraphics[width=\linewidth]{\resultdir/Ablation_E2_rank_val_loss.png}
\caption{Validation loss}
\end{subfigure}
\begin{subfigure}[b]{0.49\linewidth}
\includegraphics[width=\linewidth]{\resultdir/Ablation_E2_rank_train_loss.png}
\caption{Training loss}
\end{subfigure}
\caption{E2: validation and training loss by rank.}
\end{figure}

\begin{table}[ht]
\centering
\begin{tabular}{lc}
\toprule
\textbf{Rank} & \textbf{Test accuracy (mean $\pm$ std)} \\
\midrule
rank = 16   & $0.9068 \pm 0.0014$ \\
rank = 32   & $0.9180 \pm 0.0014$ \\
rank = 64   & $0.9284 \pm 0.0015$ \\
rank = 128  & $0.9357 \pm 0.0014$ \\
\textbf{rank = 256}  & $\mathbf{0.9381 \pm 0.0014}$ \\
\bottomrule
\end{tabular}
\caption{E2: final TTA test accuracy by rank.}
\end{table}

\subsection{E3: NS-Step Sweep}

\begin{figure}[ht]
\centering
\begin{subfigure}[b]{0.49\linewidth}
\includegraphics[width=\linewidth]{\resultdir/Ablation_E3_NS_step_val_loss.png}
\caption{Validation loss}
\end{subfigure}
\begin{subfigure}[b]{0.49\linewidth}
\includegraphics[width=\linewidth]{\resultdir/Ablation_E3_NS_step_train_loss.png}
\caption{Training loss}
\end{subfigure}
\caption{E3: validation and training loss by NS iterations.}
\end{figure}

\begin{table}[ht]
\centering
\begin{tabular}{lc}
\toprule
\textbf{NS iterations} & \textbf{Test accuracy (mean $\pm$ std)} \\
\midrule
$q = 1$  & $0.8940 \pm 0.0019$ \\
$q = 3$  & $0.9165 \pm 0.0013$ \\
$q = 5$  & $0.9296 \pm 0.0016$ \\
$q = 7$  & $0.9362 \pm 0.0012$ \\
$\mathbf{q = 9}$  & $\mathbf{0.9386 \pm 0.0015}$ \\
\bottomrule
\end{tabular}
\caption{E3: final TTA test accuracy by NS iterations.}
\end{table}

\subsection{E4: Batch-Size Sweep}

\begin{figure}[ht]
\centering
\begin{subfigure}[b]{0.49\linewidth}
\includegraphics[width=\linewidth]{\resultdir/Ablation_E4_batch_size_val_loss.png}
\caption{Validation loss}
\end{subfigure}
\begin{subfigure}[b]{0.49\linewidth}
\includegraphics[width=\linewidth]{\resultdir/Ablation_E4_batch_size_train_loss.png}
\caption{Training loss}
\end{subfigure}
\caption{E4: validation and training loss by batch size.}
\end{figure}

\begin{table}[ht]
\centering
\begin{tabular}{lc}
\toprule
\textbf{Batch size} & \textbf{Test accuracy (mean $\pm$ std)} \\
\midrule
$B = 500$   & $0.9373 \pm 0.0013$ \\
$\mathbf{B = 1000}$  & $\mathbf{0.9386 \pm 0.0016}$ \\
$B = 2000$  & $0.9339 \pm 0.0015$ \\
$B = 3000$  & $0.9272 \pm 0.0022$ \\
$B = 4000$  & $0.9120 \pm 0.0047$ \\
\bottomrule
\end{tabular}
\caption{E4: final TTA test accuracy by batch size.}
\end{table}


\end{document}
